\documentclass[journal, onecolumn]{IEEEtran}

\usepackage{titlesec}
\usepackage{graphicx}
\usepackage{array}
\usepackage{makecell}
\usepackage{float}
\usepackage{caption}
\usepackage{wrapfig}

\titleformat{\section}
    {\normalfont\bfseries}
    {\thesection}
    {1em}
    {}

\titleformat{\subsection}
    {\normalfont\bfseries}
    {\thesubsection}
    {1em}
    {}

\renewcommand\thesection{\arabic{section}}
\renewcommand\thesubsection{\thesection.\arabic{subsection}}
\renewcommand\thesubsubsection{\thesubsection.\arabic{subsubsection}}

\begin{document}

\title{FireFly Design Document}
\author{Christian Li, Kenny Wang, Kevin Yao, Matin Sobhani, Nick Lafredo, Parsa Rezvani, Ryan Yu}
\date{\today}

\maketitle

\begin{figure*}[!t]
    \centering
    \includegraphics[width=.5\linewidth]{images/UCSC_Logo.png}
    \caption*{}
    \label{fig:UCSCLOGO}
\end{figure*}

\vspace{1in}

\begin{figure*}[!t]
    \centering
    \includegraphics[width=0.25\linewidth]{images/BE_Logo.png}
    \caption*{}
    \label{fig:BELOGO}
\end{figure*}

\clearpage

\tableofcontents

\vspace{1in}

\begin{abstract}
This document presents the design and proof-of-concept for FireFly, a rapidly deployable environmental monitoring drone for wildfire detection in infrastructure-free terrain.
\end{abstract}

\newpage

% ============================================================
\section*{Executive Summary}
\addcontentsline{toc}{section}{Executive Summary}
% ============================================================

FireFly is a rapidly deployable quadcopter designed for first-response
environmental monitoring in high wildfire-risk areas with no existing sensor
infrastructure. The system collects real-time temperature, humidity, and smoke
data and transmits live readings to a ground station, triggering alerts when
fire risk thresholds are crossed.

\vspace{0.5em}

The design centers on a 3D-printed X-frame carrying a dedicated onboard microcontroller, temperature and humidity sensor, and smoke detector. A flight controller manages propulsion and stabilization while the onboard microcontroller handles sensor acquisition and wireless transmission to a ground receiver. A 3S LiPo battery powers both propulsion and electronics for the duration of a monitoring mission.

\vspace{0.5em}

Validation consists of bench testing individual modules (power, propulsion,
sensors, and communications) followed by tethered and open-field flight
testing for stable hover, data link verification, and threshold-based smoke
and temperature detection. The goal is to demonstrate that FireFly can serve
as a practical proof-of-concept for infrastructure-independent wildfire
monitoring.

\clearpage

% ============================================================
\section{Introduction}
% ============================================================

This document contains all the information needed to design and use the FireFly Drone.

\vspace{1em}

Why is collecting data during a pre-fire window important? A pre-fire window, the period before an incident occurs, is critical because it provides time-sensitive information that improves responder safety, reduces property damage, and enables faster, more effective emergency response. Data on vegetation conditions, topography, temperature, humidity, and smoke concentration helps predict fire behavior, identify high-risk areas, and support mitigation actions such as controlled burns or targeted thinning operations. Without this information, emergency crews arrive with incomplete situational awareness and must gather critical details in real time under dangerous conditions.

\vspace{1em}

For wildfire response specifically, the challenge is compounded by geography. High-risk areas are frequently roadless wilderness where permanent sensor infrastructure cannot be installed or maintained. Satellite imagery and aerial patrols provide coverage, but both suffer from response delays and gaps during the highest-risk windows. The result is a critical blind spot: the pre-fire period in remote terrain is precisely where early awareness matters most, yet it is the hardest environment in which to collect data.

\vspace{1em}

FireFly addresses this gap directly. By deploying a sensor-equipped quadcopter into areas with zero existing infrastructure, FireFly enables real-time collection of temperature, humidity, and smoke data during the pre-fire window, transmitting live readings to a ground station and triggering alerts when fire risk thresholds are crossed. The platform is designed for rapid deployment by field crews without requiring extensive training or specialized support, transforming what would otherwise be a period of uncertainty into a structured, data-driven response.

\subsection{Need Statement}
Remote and high fire-risk areas lack the sensor infrastructure necessary to monitor environmental conditions during the critical pre-fire window, leaving fire response teams without the real-time data needed to act early.

\subsection{Goal Statement}
Provide fire response teams with real-time temperature, humidity, and smoke readings from high-risk, infrastructure-free terrain by deploying a rapidly mobilized monitoring system that transmits real-time data to a ground station and alerts operators when fire risk thresholds are crossed.

\subsection{Personas}
Individuals with experience levels ranging from novice to expert in fire monitoring and environmental fieldwork.

\subsubsection*{End Users}

\begin{wrapfigure}{l}{0.2\textwidth}
    \includegraphics[width=0.15\textwidth]{images/marcus.png}
\end{wrapfigure}
\textbf{Marcus} \\
\underline{\textit{Role:}} Wildfire Operations Coordinator (professional first responder) \\
\underline{\textit{Background / Context:}} Works with a state forestry division managing prescribed burns and post-fire monitoring. Needs something fast that he can hand off to field crews without extensive training. \\
\underline{\textit{Goals:}} Rapid deployment of sensor-equipped drones into areas too dangerous or inaccessible for personnel during the high-risk pre-fire window. Wants real-time smoke and temperature data without putting people in harm's way. \\
\underline{\textit{Behaviors:}} Works under time pressure, values reliability over features, and coordinates with multiple agencies. Skeptical of unproven technology. Practical: if it breaks in the field, it needs to be fixable on-site with minimal tools. \\
\underline{\textit{Expectations:}} A rugged, simple-to-operate platform that can be deployed at a moment's notice. Wants clear threshold alerts he can act on immediately.

\vspace{1em}

\begin{wrapfigure}{l}{0.2\textwidth}
    \includegraphics[width=0.18\textwidth]{images/elena.png}
\end{wrapfigure}
\textbf{Elena} \\
\underline{\textit{Role:}} Ecology Graduate Researcher (academic user) \\
\underline{\textit{Background / Context:}} Studying fire risk and ecosystem conditions at a land-grant university. Needs to collect environmental sensor data across wilderness areas where no permanent monitoring infrastructure exists. Currently hikes in with handheld sensors, which limits how often and how broadly she can sample. \\
\underline{\textit{Goals:}} Automate pre-fire data collection over defined survey areas. Wants live temperature, humidity, and smoke readings logged without requiring physical presence at each site. \\
\underline{\textit{Behaviors:}} Comfortable with technical documentation, participates in academic drone and environmental monitoring communities, and documents everything for publication. Willing to configure and customize but does not have time to build from scratch. \\
\underline{\textit{Expectations:}} A reliable platform with a clear data pipeline she can connect to her existing research database, and results trustworthy enough to cite.

\vspace{1em}

\subsubsection*{Experts / Consultants}

\begin{wrapfigure}{l}{0.2\textwidth}
    \includegraphics[width=0.18\textwidth]{images/anita.png}
\end{wrapfigure}
\textbf{Dr. Anita Rowe} \\
\underline{\textit{Role:}} Environmental Sensor Systems Engineer \\
\underline{\textit{Background / Context:}} Works at a national lab developing atmospheric sensor arrays for wildfire monitoring. Has deep expertise in what data is actually actionable in pre-fire environments and where sensors fail under heat, humidity, and particulate exposure. \\
\underline{\textit{Goals:}} Validate that FireFly's sensor suite captures meaningful data for fire risk assessment and identify unreliable component choices before they are locked into the design. \\
\underline{\textit{Behaviors:}} Data-driven and skeptical of consumer-grade sensors in harsh field conditions. Will probe calibration drift, operating temperature limits, and whether threshold alert logic reflects real-world fire risk indicators. \\
\underline{\textit{Expectations:}} A platform sound enough in its sensor selection and alert logic that field results can be trusted and the methodology can be defended to other scientists.

\vspace{1em}

\subsubsection*{Client}

\begin{wrapfigure}{l}{0.2\textwidth}
    \includegraphics[width=0.18\textwidth]{images/usfs.png}
\end{wrapfigure}
\textbf{U.S. Forest Service Regional Office} \\
\underline{\textit{Role:}} Federal Land and Fire Management Agency \\
\underline{\textit{Background / Context:}} Responsible for fire prevention, monitoring, and response across millions of acres of national forest, much of it roadless wilderness with no existing sensor infrastructure. Currently relies on satellite imagery, fire lookout towers, and aerial patrols, all of which have coverage gaps and response delays during the highest-risk windows. \\
\underline{\textit{Goals:}} Evaluate FireFly as a rapidly deployable supplement to existing monitoring systems for areas where infrastructure cannot operate. Interested in whether the platform can reliably trigger alerts before a fire establishes. \\
\underline{\textit{Behaviors:}} Evaluates tools based on operational reliability and ease of use for non-specialist field crews. Requires demonstrated performance data before recommending adoption. \\
\underline{\textit{Expectations:}} A working prototype that proves the core concept: zero-infrastructure deployment, live sensor transmission, and threshold-based alerting, with documentation clear enough for field personnel to operate without engineering support.

\subsection{Existing Research}
Current wildfire detection relies on satellite imagery, fixed camera networks,
and manned aircraft. Satellites cannot reliably detect fires smaller than
approximately 18 square kilometers. Fixed towers like AlertCalifornia's
camera network are constrained by line of sight, leaving roadless terrain
uncovered. UAV-based systems have shown promise, with UAV--IoT networks
detecting fires as small as 2.5 square kilometers at over 99 percent accuracy,
but most depend on pre-deployed ground sensors or persistent network
infrastructure that does not exist in remote areas.

\vspace{0.5em}

FireFly addresses this gap by providing a self-contained, rapidly deployable
platform that requires no pre-existing infrastructure to deliver real-time
environmental data and threshold-based fire risk alerts.

\subsection{Sustainability Statement}

The FireFly project is committed to the ``Design for Environment'' philosophy. Our goal is to provide a high-tech solution for wildfire mitigation while ensuring that the lifecycle of the drone itself does not offset its environmental benefits.

\vspace{1em}

\textbf{Environmental Stewardship:}
The primary purpose of FireFly is the preservation of biodiversity and carbon sinks by preventing catastrophic wildfires. By providing early-warning data in the pre-fire window, we reduce the need for large-scale, carbon-intensive emergency responses such as heavy tanker aircraft and ground fleets. Furthermore, by utilizing electric propulsion over internal combustion, we ensure zero-emission operation during monitoring flights.

\vspace{1em}

\textbf{Material Selection and Circularity:}
To minimize our manufacturing footprint, we have opted for 3D-printed components using \textbf{PLA}. PLA is a bio-based, industrially compostable polymer derived from renewable resources like corn starch. By using additive manufacturing, we significantly reduce material waste compared to traditional subtractive machining.

\vspace{1em}

\textbf{Life Cycle and Battery Management:}
We recognize that Lithium-Polymer (LiPo) batteries pose a chemical risk if disposed of improperly. Our operational plan includes a strict end-of-life protocol where degraded cells are discharged and taken to certified e-waste recycling centers. We have designed the electronics with standardized connectors to ensure that sensors and flight controllers can be harvested and reused in future iterations of the platform, extending the functional life of the high-value silicon components.

\vspace{1em}

\textbf{Social and Safety Impact:}
Sustainability includes the protection of human capital. FireFly enhances the safety of first responders by providing situational awareness in hazardous terrains without risking human life. By creating a rapidly deployable tool, we extend access to advanced fire-monitoring capability to smaller forestry divisions and academic researchers who operate outside the reach of satellite-based or manned aerial surveillance programs.

\vspace{1em}

\textbf{Ethics Statement:}
As engineers, we recognize that our work has a direct impact on public safety and the environment. We are committed to upholding the highest standards of honesty, transparency, and responsibility throughout the design and development of FireFly. Our system is intended solely to support fire prevention and responder safety. We will not misrepresent the capabilities of the prototype, will clearly communicate the limitations of sensor accuracy and flight duration, and will ensure that any data collected during testing is handled responsibly. We adhere to the NSPE Code of Ethics, holding paramount the safety, health, and welfare of the public in all design decisions.

\newpage

% ============================================================
\section{Design}
% ============================================================

\subsection{Aesthetic Prototype}

The FireFly aesthetic prototype follows a field-ready design language: a compact enclosed central body with rounded corners, symmetric X-arm geometry, and cylindrical motor pods at each corner. The overall form emphasizes structural clarity and quick visual inspection rather than cosmetic styling. Rounded exterior edges reduce snag points during transport, while the open-top bay and side windows communicate where electronics, wiring, and battery hardware are intended to be serviced.

\begin{figure}[H]
    \centering
    \includegraphics[width=0.46\textwidth]{images/drone_design_5_top.png}
    \hfill
    \includegraphics[width=0.46\textwidth]{images/drone_design_5_corner.png}
    \caption{FireFly CAD aesthetic prototype, iteration v5 (top and isometric views).}
    \label{fig:drone_v5_aesthetic}
\end{figure}

The prototype evolved through multiple print-test iterations. Early versions validated basic packaging and motor placement, while later versions introduced a cleaner body-to-arm transition, larger internal volume, and improved vent/window geometry for lighter weight and faster access to internals. This iterative process preserved a consistent visual identity while converging on a more manufacturable and maintainable frame.

\begin{figure}[H]
    \centering
    \includegraphics[width=0.85\textwidth]{images/drone_design_5_compare_1.jpeg}
    \caption{Physical prototype progression (v1--v3) that informed the final v5 frame direction.}
    \label{fig:drone_proto_progression}
\end{figure}

\subsection{Design for Manufacture, Assembly, Maintenance}

\textbf{Design for Manufacture (DFM):}
The current frame is designed as a single-piece 3D print to minimize part count and remove alignment errors. The v5 CAD envelope is approximately $100 \times 72 \times 24$ mm with $2.5$ mm floor and wall thicknesses, balancing print speed, stiffness, and mass. Arms are sized at $10.0$ mm width and $5.4$ mm thickness to maintain rigidity under motor thrust and landing loads. Motor pods are dimensioned for SunnySky R1106-class motors with built-in tolerance ($19$ mm outer diameter, $15.0$ mm inner diameter, $17$ mm height, $2.2$ mm mounting holes on a $9$ mm bolt-circle).

\vspace{0.5em}

\textbf{Design for Assembly (DFA):}
Assembly features are integrated directly into the print geometry. The frame includes four battery strap slots in the floor and continuous wire-routing paths from motor pods into the electronics bay. This allows a straightforward assembly order: install motors in pods, route motor leads through arm channels, mount the controller board, then secure the battery and sensor harness. Because most interfaces are geometric rather than adhesive, rework is possible without damaging the frame.

\vspace{0.5em}

\textbf{Design for Maintenance (DFMn):}
Maintenance access is provided by the open top cavity, side-wall windows, and pod-to-arm wire notches. These features allow quick inspection of solder joints, connector seating, and wire wear without full teardown. Electronics can be removed independently from propulsion hardware, and battery straps can be replaced in the field with standard hook-and-loop straps. In case of structural damage, the one-piece architecture supports rapid replacement by reprinting the frame using the same CAD file, preserving dimensional repeatability across repairs.

\subsection{Block Diagrams}

\begin{figure}[H]
    \centering
    \includegraphics[width=1.0\textwidth]{images/Final_Schematic.png}
    \caption{A high-level schematic of the FireFly system.}
    \label{fig:schematic_design}
\end{figure}

\subsection{Wiring Diagrams}

\begin{figure}[H]
    \centering
    \includegraphics[width=0.75\textwidth]{images/Circuit Diagram.png}
    \caption{Circuit Diagram}
    \label{fig:circuit_diagram}
\end{figure}

\subsection{State Transition Diagrams}

\begin{figure}[H]
    \centering
    \includegraphics[width=0.75\textwidth]{images/State Transition Diagram.png}
    \caption{State Transition Diagram}
    \label{fig:state_transition}
\end{figure}

\subsection{Technology (Languages, Hardware, etc.)}

FireFly uses a two-node wireless architecture connected via a low-latency peer-to-peer radio protocol. The drone carries an onboard microcontroller that reads sensor data from a temperature and humidity sensor and a smoke detector, packages the readings into a structured telemetry packet, and transmits it wirelessly to a ground receiver. The communication link operates on the 2.4GHz band with a range of approximately 400 meters line-of-sight and latency under 10ms, making it well suited for remote wildfire monitoring scenarios where reliable low-latency telemetry is critical.

\vspace{0.5em}

The ground station consists of a second microcontroller that acts as a dedicated receiver, unpacks the incoming telemetry data, and forwards it as a JSON stream over USB serial to a host machine. The host machine runs a Python Flask server that reads the serial stream and exposes a local HTTP endpoint, which a browser-based dashboard polls at one-second intervals to display live sensor readings. Firmware on both microcontrollers is written in C, with the ground station dashboard implemented in Python and JavaScript.

\subsection{Simulations}

Before building and flying, we used simulation tools to verify that the FireFly configuration can sustain flight, carry its sensor payload, and meet the core mission requirement of delivering real-time environmental data during the pre-fire window. An online multicopter calculator (eCalc xcopterCalc) was used to evaluate combinations of frame size, motors, propellers, and battery against thrust-to-weight ratio, hover throttle, current draw, and expected flight time under sensor payload assumptions. This confirmed that the selected hardware provides sufficient thrust margin to carry the temperature/humidity sensor and smoke detector while maintaining acceptable hover efficiency and a flight time long enough to collect and transmit meaningful readings to the ground station.

\vspace{0.5em}

A simplified quadcopter dynamics model was used to tune attitude and altitude control parameters, verifying that the platform can hold a stable position over a monitoring area, a prerequisite for consistent sensor sampling. Research on quadcopter modeling confirms that such simulations can predict controller behavior and reduce the number of risky test flights needed to achieve stable flight. Together, these simulations provide evidence that FireFly can meet its need statement requirements: reaching areas with zero sensor infrastructure and sustaining flight long enough to collect and transmit temperature, humidity, and smoke data before committing to full implementation and field testing.

\newpage

% ============================================================
\section{Evaluation}
% ============================================================

\subsection{Functional Prototype}
As of the initial submission, FireFly's functional prototype is currently under assembly. Core components, including the onboard microcontrollers, temperature and humidity sensor, and smoke detector, have been sourced and bench-tested for communication and sensor output. A prototype circuitry system using an H-Bridge in place of the ESC and a DC brushed motor in place of the brushless motor was assembled to validate the wireless communication link between the ground-side and drone-side microcontrollers while awaiting final parts. Full integration with the flight stack and ground station dashboard is in progress. Photographs and detailed performance data will be included in the final Spring Quarter submission.

\subsection{Testing}
Initial bench testing has validated point-to-point wireless communication between the drone-side and ground-side microcontrollers, confirmed sensor readout from the temperature/humidity and smoke sensors, and verified JSON telemetry forwarding over USB serial to the Flask dashboard. Full tethered and open-field flight testing with threshold-based alert verification is planned for the prototype phase. Detailed test procedures and results are documented in Appendix 3.

\newpage

% ============================================================
\section{Parts List}
% ============================================================

\begin{table}[H]
    \centering
    \small
    \begin{tabular}{|l|c|c|}
        \hline
        \textbf{Part} & \textbf{Supplier} & \textbf{Quantity} \\
        \hline
        CarbonX PLA+CF Filament & Creality & 1 \\
        \hline
        Battery: LiPo 850 mAh 3S 25/30C & HobbyKing & 1 \\
        \hline
        Motor: SunnySky R1106 FPV Brushless Motors & Druav & 4 \\
        \hline
        Propeller: Gemfan Hurricane 2009-3 Tri-Blade (Set of 4) & Pyrodrone & 1 \\
        \hline
        HAKRC 20A 4-in-1 \verb|3-6S BLHeli_32| 5V 3A BEC DShot1200 Mini ESC & Alibaba & 1 \\
        \hline
        Perforated Circuit Board & Adafruit & 2 \\
        \hline
        Female Header Pins 2x6 & AliExpress & 4 \\
        \hline
        Female Header Pins 2x8 & AliExpress & 2 \\
        \hline
        Female Header Pins 2x10 & Mikroe & 2 \\
        \hline
    \end{tabular}
    \caption{Itemized Parts List}
    \label{tab:budget}
\end{table}

\newpage

% ============================================================
\section{Appendix 1 - Problem Formulation}
% ============================================================

\subsection{Conceptualisations}

FireFly is a lightweight quadcopter designed to collect and transmit
environmental sensor data in remote areas with no existing infrastructure.
The system is built around three functions: stable flight, environmental
sensing, and real-time data transmission.

\vspace{0.5em}

The drone carries an onboard microcontroller paired with a temperature and humidity sensor and a smoke detector. A flight controller manages propulsion and stabilization while the onboard microcontroller handles sensor acquisition and wireless transmission. A LiPo battery powers both propulsion and electronics for the duration of a monitoring mission.

\vspace{0.5em}

Data is transmitted wirelessly to a ground receiver microcontroller, which forwards readings over USB serial to a host machine running a Python Flask dashboard displaying live sensor data and threshold alerts.

\subsection{Brainstorming}
Several frame geometries were evaluated including X and H-style quadcopter
layouts. An X-frame was selected for its balanced motor distribution,
stability, and ease of sensor mounting.

\vspace{0.5em}

The sensing approach was inspired by existing ground-based forest sensor
networks, which are geographically limited and difficult to scale into remote
terrain. FireFly was conceived as a mobile alternative: a drone that could
collect the same environmental data in areas where fixed infrastructure does
not exist.

\vspace{0.5em}

For communication, Bluetooth Low Energy and Wi-Fi were both evaluated. A connectionless peer-to-peer wireless protocol was ultimately selected for its simplicity, low latency, and sufficient range for a point-to-point link between the drone and ground receiver without the overhead of a full Wi-Fi stack.

\vspace{0.5em}

Power and propulsion options were compared using eCalc to evaluate
thrust-to-weight ratio, hover current, and projected flight time across
different motor, propeller, and LiPo configurations. This guided selection
of a 3S LiPo and motor pairing that balances flight endurance with the
electrical load of the sensor and communication hardware.

\subsection{Decision Tables}

The following decision tables evaluate key design alternatives across weighted criteria. Scores are rated 1--5 per criterion, multiplied by the criterion weight, and summed to produce a weighted total.

\begin{table}[H]
\centering
\small
\begin{tabular}{|l|c|c|c|c|}
\hline
\textbf{Criterion} & \textbf{Weight} & \textbf{X-Frame} & \textbf{H-Frame} & \textbf{+ Frame} \\
\hline
Stability & 25\% & 5 & 3 & 4 \\
\hline
Sensor Mounting & 20\% & 5 & 4 & 3 \\
\hline
Ease of Assembly & 20\% & 4 & 4 & 3 \\
\hline
Printability & 20\% & 4 & 3 & 5 \\
\hline
Repairability & 15\% & 4 & 5 & 3 \\
\hline
\textbf{Weighted Total} & & \textbf{4.45} & \textbf{3.75} & \textbf{3.65} \\
\hline
\end{tabular}
\caption{Frame Geometry Decision Table}
\end{table}

\begin{table}[H]
\centering
\small
\begin{tabular}{|l|c|c|c|c|}
\hline
\textbf{Criterion} & \textbf{Weight} & \textbf{Peer-to-Peer Wireless} & \textbf{Wi-Fi} & \textbf{BLE} \\
\hline
Latency & 30\% & 5 & 3 & 4 \\
\hline
Range & 25\% & 5 & 4 & 2 \\
\hline
Infrastructure Required & 25\% & 5 & 2 & 4 \\
\hline
Implementation Complexity & 20\% & 4 & 3 & 4 \\
\hline
\textbf{Weighted Total} & & \textbf{4.80} & \textbf{3.00} & \textbf{3.50} \\
\hline
\end{tabular}
\caption{Communication Protocol Decision Table}
\end{table}

\begin{table}[H]
\centering
\small
\begin{tabular}{|l|c|c|c|}
\hline
\textbf{Criterion} & \textbf{Weight} & \textbf{PLA} & \textbf{Carbon Fiber PLA} \\
\hline
Weight & 30\% & 4 & 5 \\
\hline
Printability & 25\% & 5 & 3 \\
\hline
Structural Rigidity & 25\% & 3 & 5 \\
\hline
Repairability & 20\% & 5 & 3 \\
\hline
\textbf{Weighted Total} & & \textbf{4.20} & \textbf{4.00} \\
\hline
\end{tabular}
\caption{Frame Material Decision Table}
\end{table}

\subsection{Morphological Charts}
The following combinations were explored for the primary subsystems:
\begin{figure}[H]
    \centering
    \includegraphics[width=0.5\linewidth, height=0.45\textheight]{images/morphological.png}
    \caption{Morphological analysis chart outlining potential components of FireFly}
    \label{fig:morphological}
\end{figure}

\newpage

% ============================================================
\section{Appendix 2 - Planning}
% ============================================================

\subsection{Basic Plan / Gantt Chart}
\begin{figure}[H]
    \centering
    \includegraphics[width=1\linewidth, height=0.50\textheight, keepaspectratio]{images/Gantt_Chart_v2.png}
    \caption{Gantt Chart}
    \label{fig:gantt_chart}
\end{figure}

\subsection{Division of Labor}

\begin{table}[H]
    \centering
    \begin{tabular}{|l|l|}
        \hline
        \textbf{Team Member} & \textbf{Primary Responsibility} \\
        \hline
        Christian Li & Flight Controller Programming \\
        \hline
        Kenny Wang & Hardware Integration \\
        \hline
        Kevin Yao & Mechanical Design / CAD \\
        \hline
        Matin Sobhani & Mechanical Design / CAD \\
        \hline
        Nick Lafredo & Power Systems / Documentation \\
        \hline
        Parsa Rezvani & Testing and Validation / Telemetry \\
        \hline
        Ryan Yu & Flight Controller Programming \\
        \hline
    \end{tabular}
    \caption{Division of Labor during Prototyping}
\end{table}

\subsection{Collaboration}
\begin{itemize}
    \item \textbf{Github:} CI/CD Tool
    \item \textbf{Discord:} Communication and meetings
    \item \textbf{Notion:} Gantt Chart
    \item \textbf{Trello:} Kanban board for task tracking
\end{itemize}

\newpage

% ============================================================
\section{Appendix 3 - Test Plan \& Results}
% ============================================================

The test plan consists of three phases:
\begin{enumerate}
    \item \textbf{Bench Testing:}

    We began by creating a prototype circuitry system that utilizes an H-Bridge in substitution for the Electronic Speed Controller, along with a DC brushed motor in place of the designed DC brushless motor. This system was used to test the wireless communication between the ground-side and drone-side microcontrollers, allowing us to develop and validate the data transmission pipeline while awaiting final parts.

    \item \textbf{Tethered Flight:}

    We will integrate the ESC and DC brushless motors into the physical 3D-printed frame and test flight capability in a controlled indoor environment. During this phase we will emulate conditions that would flag a fire danger zone (elevated temperature, low humidity, and elevated smoke readings) and verify that threshold alerts are correctly triggered and transmitted to the ground station.

    \item \textbf{Open Field Flight:}

    We will test the drone in an open field to evaluate durability, structural integrity, and real-world sensor performance. Outdoor conditions including wind and heat will allow us to assess how environmental factors affect flight stability and sensor accuracy, providing direct feedback on both the mechanical design and the telemetry pipeline under realistic deployment conditions.
\end{enumerate}

\newpage

% ============================================================
\section{Appendix 4 - Review}
% ============================================================

\textit{Reflection of each team member's thoughts on the project.}

\vspace{1em}

\textbf{Christian Li:} Looking back on this project, one of the biggest lessons I learned is the importance of starting with a clearly defined need rather than a preconceived solution. We approached the project wanting to build a drone, and only later did we step back and ask why a drone was the right solution for the problem. If we had focused earlier on the lack of fire risk data in remote areas, we could have identified the required characteristics sooner and avoided some of the design and task adjustments we made midway through the quarter. Even so, our team communicated well and adapted quickly. The flexibility helped us stay on track and produce a coherent design. This experience also gave me valuable insight into working as a team, improving how we communicate, distribute work, and learn new tools and environments together.

\vspace{1em}

\textbf{Kenny Wang:} Restarting this project I think that we should focus on not building a product that interests us but develop the idea based on a universal problem or a need. From there we can then focus on implementing it in a way that can handle this universal problem or need. For example, if we had started with the problem being that we have no way of reaching remote areas with no fire detection, we can then build upon that and not just specifically building a flying drone but maybe a rover, or even something that does not even have to move. There would be a variety of implementations of helping alleviate the same problem and need while having numerous amounts of angles to implement that solution. Even with some setbacks that we faced during the beginning stages of the project we were able to have good communication within the team to implement these changes so that we were able to complete our design work.

\vspace{1em}

\textbf{Kevin Yao:} Reflecting back to the beginning of this project, I think we should have focused on the process instead of the end goal of the design we wanted to make. One thing that we did well was forming our group and communicating often. We were vocal on what we wanted to do and had ideas on what the design would look like. A problem that occurred was our need and goal statement. What we should have done differently was to define clear objectives of what needed to be done and check the requirements. Another thing we should have done earlier was setting up the github with our other tools to have a cleaner pipeline of our progress. Although we had some setbacks, our team was able to implement these changes quickly and improved our communication and workflow.

\vspace{1em}

\textbf{Matin Sobhani:} If we were starting this project again, I would make sure we had a solid need/goal statement, and verified it with the TA and professor. We thought we had a solid one early on around week 4, but ended up having to change it around week 8. Thankfully this did not affect us too negatively since our underlying prototype idea stayed mostly the same, but it did cost us time and forced us to reframe parts of our documentation. In hindsight, we should have scheduled a formal requirement check-in earlier and created clearer milestones tied to those requirements so we could catch misalignment sooner. What went well is that we stayed very communicative throughout the project, kept each other up-to-date, and adapted quickly whenever priorities changed.

\vspace{1em}

\textbf{Nick Lafredo:} This project is the first I've done on such a large scale, and it taught me a lot about the timeline and organization needed to accomplish a group goal. If we were to go back and start a new project, we would focus more on the need and solution, rather than starting with a product we wanted to make and finding a need that fits it. It is also now much clearer that we are not designing the prototype, but creating a general design that could be implemented in the future, with the prototype proving its possibility. We faced a few setbacks with our design and deciding what problem we were aiming to solve, but ultimately, it made our final idea much stronger and well-defined. Going forward, we will have more frequent check ins with our team and our instructors to make sure we stay on the right track. We now have a clear path towards completing the design, and I am confident that we will reach our goal.

\vspace{1em}

\textbf{Parsa Rezvani:} This project has been a significant learning experience for both me and the team. Initially, we approached it as an exercise in selecting an interesting idea and producing documentation around it. Over time, however, we came to understand that the class demanded something closer to real-world engineering practice: defining a defensible need, building toward a clear goal, and maintaining a concept that could stand on its own merits in a professional context. Had we internalized that framing earlier, our process would have been more focused from the start. Despite the adjustments along the way, I am proud of how the team collaborated and adapted. With our components arriving soon, I am confident we are well positioned to make meaningful progress on the prototype.

\vspace{1em}

\textbf{Ryan Yu:} Looking back at the beginning of this project, there are a few things that went well and a few things that went poorly. I think we got our group together, talked about initial parts, calculations, and ideas quite early, and were on top of things. However, I think the communication between the professor, TA, and our group was extremely poor. It was not clearly laid out that we were designing a product but producing a prototype. We had difficulties having useful long discussions with the professor and were often confused on what he wanted, constantly changing last minute.

\end{document}
